\section{Threats to validity}
\label{research-design-and-methodology:threats-to-validity}

Construct validity concerns the gap between what an instrument records and what the experiment intends to measure, and two outcomes are exposed to it. The first is network transfer, counted at the host boundary (Section~\ref{research-design-and-methodology:metrics-and-cost-model:primary-metrics}). Transport framing, encryption, and retransmission below that boundary are therefore not separated out, so the metric approximates the bytes moved over the network rather than measuring them exactly. Because the same boundary applies to every configuration, this bears on absolute figures more than on comparisons. The k-nearest-neighbor workload raises the second concern. Distance ordering is index-served only with a bounding-box predicate (Section~\ref{background:spatial-data-formats-and-storage-models:full-database-engines}), so engines may rank by a bounding-box proxy before refining to exact distance. Fixing the query point and the returned count bounds this risk but does not remove it: two engines may still order candidates by different internal distances. Distinct from these two instrument-level outcomes, the single-node run-time comparison of \acrshort{rq}1 carries a construct concern at the level of the comparison: because each engine ran in its native configuration, the run-time ordering measures spatial-indexing capability and storage locality together with the engine rather than the engine in isolation, with PostGIS both indexed and co-located and the GeoParquet path neither. That bundling is the configuration factor under test by design, so the ordering is valid for the stacks as deployed, but it bounds attribution to the engine alone, the concrete per-engine configurations being given in Section~\ref{system-architecture-and-implementation:compute-configuration}.

The sequential stopping rule in Section~\ref{research-design-and-methodology:workloads:iteration-counts} raises a further construct concern. The rule halts a single-machine workload once the bootstrap confidence interval on its mean elapsed time reaches the target half-width, but never before a \num{60}-second timed window has elapsed. That floor follows from the metric-collection granularity the cost model relies on (Section~\ref{system-architecture-and-implementation:measurement-instrumentation-and-cost-model:host-level-instrumentation}). For a sub-millisecond query it dominates the budget. Most of the wall-clock time goes to satisfying the window rather than narrowing the interval, so the achieved iteration count reflects the floor more than the variance it was meant to track. The interval is itself noisy when it rests on fewer than \num{30} samples, and the \num{10}-iteration minimum that every workload must clear bounds this noise without removing it.

Internal validity asks whether an observed difference reflects the factor under test rather than a confound, and the dominant confound here is the variability of shared cloud infrastructure. Production cloud services vary over time and between nominally identical allocations \parencite{iosup2011_performance_variability_cloud_services}, and instances of the same specification differ measurably in processor, input/output, and network throughput \parencite{schad2010_runtime_measurements_cloud}. The four measures of Section~\ref{research-design-and-methodology:overview-and-experimental-design:replication-and-randomization} constrain this confound rather than assume it away: concurrent pairing on identically provisioned instances, seeded randomization of execution order, warmup to steady state, and a fresh cluster per distributed run. Running test and control on one instance in randomized order is what lets a small difference register as an effect instead of as drift. When a batch member fails mid-run, for instance through cluster-provisioning quota exhaustion, the orchestrator detects the failure and records whether it was partial or total, so the affected batches can be identified during analysis. A residual threat remains: a surviving peer then ran in the same wall-clock window without the contention its paired counterpart would have introduced, which may bias its measurement favorably.

A further concern is that the framework distinguishes the two failure modes only after the fact, not at runtime. Deterministic resource-exhaustion failures recur on every attempt and signal that the configuration cannot run the workload at that scale, as when the distributed engine runs out of memory at the smallest cluster size on the largest tier, or the single-node relational engine runs out of memory on the national-scale join. Transient cluster flakiness is incidental, such as a lost-executor failure on a large-tier run that succeeds on retry. The framework tolerates the transient case, recording each failed iteration as a sample marked \texttt{failed} without aborting the run. While a run is in progress, however, it cannot separate the two modes. That distinction falls to downstream analysis through the failed-iteration count and the per-iteration failure reason. A run carrying the stop reason \texttt{partial} is retained in the analysis, but its sample is smaller than intended and may be drawn from a non-representative subset, because whatever a failed early iteration would have measured is absent from the distribution.

External validity concerns how far the findings extend beyond their conditions. The design runs in a single region, with a single provider, on one dataset domain. The architectural orderings are expected to carry beyond these conditions: which access pattern moves fewer bytes, and whether distribution repays its coordination cost. They follow from format and engine structure (Section~\ref{background:spatial-data-formats-and-storage-models:summary}; Section~\ref{background:cloud-native-data-architecture:distributed-and-non-distributed-computation}) rather than from the platform, which is why they carry. Absolute times and costs do not carry, because they depend on provider- and region-specific rates whose predictability itself varies \parencite{leitner2016_patterns_chaos_iaas}. Data is the most limited axis. A workload governed by spatial skew (Section~\ref{background:distributed-spatial-processing:spatial-partitioning-and-parallel-efficiency}), such as the distributed join, is the most sensitive to building polygons not representing other inputs. Table~\ref{tab:threats-validity-summary} collects these threats across all three validity dimensions, with the design measures that address them.

\begin{longtable}{@{}l p{0.40\textwidth} p{0.40\textwidth}@{}}
	\toprule
	Validity & Threat & Mitigation \\
	\midrule
	\endfirsthead
	\multicolumn{3}{@{}l}{\small\itshape Table~\thetable\ (continued)} \\
	\toprule
	Validity & Threat & Mitigation \\
	\midrule
	\endhead
	\midrule
	\multicolumn{3}{r@{}}{\small\itshape Continued on next page} \\
	\endfoot
	\bottomrule
    \addlinespace[\abovecaptionskip]
	\caption[Threats to validity and their mitigations]{Threats to validity grouped by dimension, the threat each poses, and how the design addresses it.}
	\label{tab:threats-validity-summary} \\
	\endlastfoot
	Construct & Bytes counted at the host boundary; sub-boundary framing, encryption, and retransmission not isolated. & Same boundary for every configuration, so it affects absolute figures more than comparisons. \\
	\addlinespace
	Construct & $k$-nearest-neighbor distance ordering is index-served only with a bounding-box predicate, so engines may use different internal distances. & Fixed query point and returned count bound the effect without removing it. \\
	\addlinespace
	Construct & The \num{60}-second timed-window floor trades the cost model's metric-collection granularity against stopping efficiency, so a sub-millisecond query may exhaust most of its budget inside the floor rather than narrowing its confidence interval. & The floor is an accepted compromise; affected workloads are identifiable from their achieved iteration count, and the driving granularity is given in Section~\ref{system-architecture-and-implementation:measurement-instrumentation-and-cost-model:host-level-instrumentation}. \\	\addlinespace
	Construct & A bootstrap confidence interval computed from fewer than \num{30} samples is noisy. & The \num{10}-iteration minimum mitigates but does not remove this. \\
	\addlinespace
	Construct & The RQ1 Shapefile configuration is exercised through GeoPandas rather than DuckDB, so its results conflate engine and format. & Acknowledged as a limitation; the configuration is read as the GeoPandas and Shapefile stack, not as DuckDB and Shapefile. \\
	\addlinespace
	Construct & Each RQ1 engine runs in its native configuration, so the single-node run-time ordering varies spatial-indexing capability and storage locality together with the engine: PostGIS is indexed and co-located, the GeoParquet path neither. & The configuration is the factor under test, so the ordering is read for the stacks as deployed, not as an isolation of the engine; the per-engine setup is given in Section~\ref{system-architecture-and-implementation:compute-configuration}. \\
	\addlinespace
	Construct & The bounding-box workload applies an area-filter predicate beyond \texttt{ST\_Intersects}, which changes selectivity across tiers and adds a reprojection cost. & Documented in the workload definition; the area cutoff is constant across configurations, so the relative comparison is preserved. \\
	\addlinespace
    Construct & The distributed engine's runtime can silently bypass the intended spatial join strategy and fall back to a brute-force cross-product join, changing the operation measured without raising an error. & The fallback path is disabled and the physical plan is verified to confirm the intended spatial join; the concrete configuration is given in Section~\ref{system-architecture-and-implementation:compute-configuration:apache-sedona-on-databricks}. \\
    \addlinespace
	Construct & The relational engine's planner statistics are stale immediately after a bulk load, which can mislead the query planner. & Statistics are refreshed after every load and a spatial index is built on both sides of the join; the concrete commands are given in Section~\ref{system-architecture-and-implementation:data-layer:postgis-on-azure-postgresql}. \\
    \addlinespace
	Internal & Shared cloud infrastructure varies over time and across nominally identical instances. & Concurrent pairing, seeded order randomization, warmup, and a fresh cluster per distributed run as described in Section~\ref{research-design-and-methodology:overview-and-experimental-design:replication-and-randomization}. \\
	\addlinespace
	External & A single region, provider, and dataset domain; the data domain is the most limited axis. & Architectural orderings are expected to carry; absolute times and costs do not. \\
	\addlinespace
    External & Cross-region egress is priced from a fixed per-gigabyte rate that may grow stale as cloud pricing changes. & The rate is pinned with the unit prices of Section~\ref{system-architecture-and-implementation:measurement-instrumentation-and-cost-model:cost-model-implementation} and flagged for revision if pricing changes. \\
\end{longtable}

The study still accepts the Shapefile read asymmetry by design and records it against the comparison it affects. Shapefile is read from local disk (Figure~\ref{fig:design-space}) and has no cloud-native read path, which favors it on transfer, so a Shapefile transfer comparison reads as a lower bound rather than a like-for-like figure. The cross-region transfer that only the distributed configuration incurs was previously treated as an accepted asymmetry, but it is now measured and priced as an explicit egress term, no longer left uncosted (Section~\ref{research-design-and-methodology:metrics-and-cost-model:cost-model}). The warmup that the distributed runs perform against their reused cluster removes the cold-start asymmetry earlier attributed to them (Section~\ref{research-design-and-methodology:overview-and-experimental-design:replication-and-randomization}). The remaining Shapefile asymmetry is not hidden in an aggregate. It is reported against the workload and dimension it touches.